\section{Extension to richer pair representations}
\label{sec:richer-relations}

The main study uses only the four standard ParT pair features. Here we
extend that representation with three physics-inspired relation families
and examine whether the EdgeValue gains persist. This supplementary study
varies the pair representation while retaining the attention and value-message
mechanisms described in Section~\ref{sec:method}.

\subsection{Additional relation descriptors}
\label{sec:selected-descriptors}

For the comparison in this appendix, we denote the main study's four-feature
pair input $r_{ij}$ by $r^{\mathrm{std}}_{ij}$.
The selected input concatenates these standard channels with learned PT,
TRACK, and REGION embeddings of dimensions $8$, $12$, and $12$, respectively,
giving $r^{\mathrm{sel}}_{ij}\in\mathbb R^{36}$. The common width-64 stem now
receives $r^{\mathrm{sel}}_{ij}$ instead of $r^{\mathrm{std}}_{ij}$.
Throughout, $i$ denotes the query and $j$ the context, and directional features
are not symmetrized. TRACK and REGION use pseudorapidity for angular
separations, whereas the standard ParT pair features use rapidity.

\paragraph{PT.}
Let $S_T=\sum_{k\in\mathcal V}p_{T,k}$ be the scalar constituent-$p_T$ sum
and $\epsilon=10^{-6}$ the numerical safeguard used by the added relation
families. Each endpoint supplies $p_{T,i}/(S_T+\epsilon)$, the safeguarded
log-fraction $\log[(p_{T,i}+\epsilon)/(S_T+\epsilon)]$, and its descending
normalized $p_T$ rank. The remaining four channels are the context-minus-query
log-fraction and rank differences, the pair log-fraction
$\log[(p_{T,i}+p_{T,j})/(S_T+\epsilon)+\epsilon]$, and
$(p_{T,j}-p_{T,i})/(p_{T,j}+p_{T,i}+\epsilon)$. Tied ranks are averaged.
A $10\!\to\!32\!\to\!8$ network embeds these ten quantities.

\paragraph{TRACK.}
Each endpoint supplies $d_0$, $d_z$, logarithmic effective uncertainties,
inverse-hyperbolic-sine impact-parameter significances, and a validity flag.
Validity requires charged-hadron, electron, or muon identification, finite,
nonmissing measurements, and positive uncertainties. Effective errors satisfy
$\sigma_{u,\mathrm{eff}}^2=\sigma_u^2+f_u^2$ for $u\in\{d_0,d_z\}$, with
floors $f_u$ fitted on training data. Shared endpoint weights produce
$g_i,g_j\in\mathbb R^{16}$ using a $7\!\to\!32\!\to\!16$ network.
The pair encoder receives
\begin{equation}
 [g_i,g_j,g_j-g_i,g_i\odot g_j,s_{ij},c_{ij}]\in\mathbb R^{81},
 \label{eq:track-pairs}
\end{equation}
where $s_{ij}$ is a four-component one-hot encoding of the ordered endpoint
validity states and $c_{ij}$ contains 13 compatibility channels. These comprise
uncertainty-normalized context-minus-query displacement differences, their
squared-sum $\chi^2$ through $\log(1+\chi^2)$ and
$\exp[-\min(\chi^2,25)/2]$, extrema of absolute transformed significances and
their signed products, and angular descriptors $\sin\Delta\phi$,
$\cos\Delta\phi$, and $\log(\Delta R+\epsilon)$, with
$\Delta\phi=\phi_i-\phi_j$ wrapped periodically and $\Delta R$ measured in
pseudorapidity--azimuth space. Compatibility channels vanish unless both
tracks are valid. The resulting $81\!\to\!48\!\to\!12$ network retains
missing-measurement information through the validity states.

\paragraph{REGION.}
A deterministic, beam-free angular tree repeatedly merges the closest pair
in $(\eta,\phi)$ using four-vector addition. It follows the angular-clustering
motivation of Cambridge/Aachen~\cite{Dokshitzer1997Cambridge,Wobisch1999Cambridge},
with the explicit pseudorapidity-based prescription above. Exclusive partitions
at $K=2,4,8$ clusters, capped by constituent multiplicity, provide 41 features:
three same-cluster indicators; five lowest-common-ancestor descriptors
(normalized depth, merge angle, $k_t$, momentum sharing, and mass); 18 endpoint
cluster $p_T$, mass, and multiplicity fractions; six constituent shares within
those clusters; six constituent-to-cluster-axis distances; and three
context-minus-query cluster-$p_T$ rank differences. Angles use reference radius
$R_{\mathrm{ref}}=0.8$, merge $k_t$ uses $p_{T,\mathrm{jet}}R_{\mathrm{ref}}$,
and mass uses jet mass. Cluster $p_T$ and mass fractions, and merge angle,
$k_t$, and mass descriptors use logarithms. Undefined diagonal merge quantities
are zeroed. A $41\!\to\!32\!\to\!12$ network embeds the result. These descriptors
are supplied for every valid pair, while attention remains global.

Here jet $p_T$ and mass are obtained from the summed constituent four-vector,
not the scalar sum $S_T$ used by PT. Ancestor depth is divided by the maximum
leaf depth, with the denominator bounded below by one.

\subsection{Preprocessing and experimental design}
\label{sec:richer-setup}

All family networks use GELU and RMS normalization between their two linear
maps and are trained jointly with the classifier. Before entering these
networks, designated continuous features are centered and rescaled using
statistics from a fixed subset of 50,000 training jets. This subset is used
only to estimate preprocessing constants; the classifier is trained on the
full one-million-jet training sample. For each feature, we subtract its median,
divide by a robust measure of its spread, and clip extreme values:
\begin{equation}
 \widehat x=\operatorname{clip}_{[-8,8]}
 \left(\frac{x-\operatorname{median}(x)}
 {\max[(Q_{75}-Q_{25})/1.349,10^{-6}]}\right).
 \label{eq:relation-normalization}
\end{equation}
Here $Q_{25}$ and $Q_{75}$ are the training-sample quartiles; the resulting
scale is less sensitive to outliers than a standard deviation. Track
uncertainty floors are also estimated from training data. Binary and
designated fixed-scale features are left unchanged. These preprocessing
constants remain fixed during validation and testing and are shared across
the selected-input configurations. Padded pair embeddings are zeroed.

The supplementary study repeats the shared/layerwise bias and EdgeValue
off/on comparisons with the richer representation, giving four additional
configurations. For the appendix comparisons, \texttt{Std} denotes the
standard pair inputs used in the main study and \texttt{Sel} denotes the
richer input. The following \texttt{S} or \texttt{L} denotes shared or
layerwise bias, and \texttt{+EV} indicates EdgeValue. Thus, \texttt{Std-S}
is the main-study ParT baseline, while \texttt{Sel-L+EV} uses the added
relations, layerwise bias, and EdgeValue.
All four additional configurations use the same data splits, three matched random seeds,
training protocol, and checkpoint-selection rule as the main study.
The selected relations use the same constituents and track measurements,
without additional detector information. Their encoders add trainable
parameters, so comparisons with standard inputs are not capacity matched.
The complete eight-configuration results appear in
Appendix~\ref{sec:complete-results}.

\subsection{Results with richer pair representations}
\label{sec:results-relations}

Table~\ref{tab:overview} includes the four selected-input configurations.
At fixed bias scheme and EdgeValue setting, adding these relations improves
mean accuracy in every comparison. The gains are $0.400$ pp for shared bias
without EdgeValue and $0.130$ pp for shared bias with EdgeValue. For layerwise
bias, the corresponding gains are $0.654$ and $0.409$ pp. All four
selected-versus-standard comparisons improve accuracy in each of the three
matched seeds.

The bias-scheme comparison also changes with the richer relation input:
the selected-relation layerwise configuration exceeds its shared counterpart
by $0.294$ pp without EdgeValue and $0.243$ pp with EdgeValue, consistently
across seeds. Additional relation information thus has a larger measured
benefit in the layerwise implementations. This is an empirical comparison of
the complete configurations, subject to the trimming and pair-processing
differences noted in Section~\ref{sec:training}, rather than a causal
attribution to separate layerwise bias maps.

EdgeValue remains beneficial after adding the selected relations, contributing
$0.385$ pp with shared bias and $0.334$ pp with layerwise bias, again with
positive gains in all three seeds. The improvements are not simply additive:
the marginal EdgeValue gain is smaller with the selected features than with
the standard features. Nevertheless, the selected-relation layerwise
ParT-EV achieves the highest mean accuracy and macro AUROC among the eight
configurations: $83.847\pm0.042\%$ and $98.399\pm0.014\%$, respectively.
Relative to baseline, these are improvements of $1.028$ pp in accuracy and
$0.170$ pp in macro AUROC. These total gains combine the change in pair
representation, bias scheme, and value messages; they do not measure
EdgeValue alone.

Class-wise gains are not uniform (Table~\ref{tab:rejection-benchmark}).
Relative to baseline, selected-relation layerwise ParT-EV increases mean
rejection at 50\% signal efficiency by $28.3\%$ for $H\to c\bar c$,
$30.8\%$ for $H\to4q$, and $37.1\%$ for fully hadronic top jets.
At the leptonic benchmark efficiencies, the gains are $23.5\%$ for
$H\to q\bar q'\ell\nu$ and $21.9\%$ for $t\to b\ell\nu$.
The matched comparison with the selected-relation layerwise model without
EdgeValue improves mean rejection in eight of the nine classes.
The exception is $H\to b\bar b$: rejection is $5253$, compared with $5271$
for baseline, and below the standard shared-bias EdgeValue model. The small
$-0.3\%$ difference from baseline does not by itself establish a degradation.
At 30\% efficiency (Table~\ref{tab:rejection-30}), gains over baseline
include $42.9\%$ for $H\to4q$ and $56.2\%$ for fully hadronic top jets.

These results show that EdgeValue can also exploit a richer pair
representation, but do not identify an optimal set of relations. PT, TRACK,
and REGION are added together, so their individual contributions and
interactions remain unresolved. A broader study of which physics-inspired
relations most benefit ParT-EV, including individual-family ablations and
capacity-matched controls, is an important direction for future work.
